%% ma: L10-B14-0085
%% lop: 10
%% chuong: 5
%% bai: 14
%% dang: 10.14.12
%% loai: DS
%% muc: [TH, NB, TH, VD]
%% dap_an: "SĐĐS"
%% nguon: "chuyen-de-cac-so-dac-trung-cua-mau-so-lieu-khong-ghep-nhom-toan-10-chuong-trinh-moi.pdf (GV. Phan Nhật Linh) (P05), Bài 03 Dạng 2 Phần II Câu 2"
%% kien_thuc: "Lớp 10 Bài 13 (số trung bình, trung vị, tứ phân vị, mốt), Bài 14 (khoảng biến thiên, khoảng tứ phân vị, phương sai, độ lệch chuẩn, giá trị bất thường)"
%% trang_thai: chinh-thuc
%% ghi_chu: "Tự duyệt theo yêu cầu giáo viên 22/10/2026."
\begin{ex}
Mẫu số liệu sau cho biết chiều cao của $10$ học sinh (đơn vị: cm):
\begin{center}
\begin{tabular}{|c|c|c|c|c|c|c|c|c|c|}\hline
$165$ & $155$ & $160$ & $145$ & $157$ & $162$ & $148$ & $170$ & $172$ & $152$\\\hline
\end{tabular}
\end{center}
Xét tính đúng sai của các khẳng định sau:
\choiceTF{Chiều cao trung bình của $10$ học sinh là $157{,}6$}{\True Khoảng biến thiên của mẫu số liệu là $27$}{\True Khoảng tứ phân vị của mẫu số liệu là $13$}{Phương sai của mẫu số liệu nhỏ hơn $64$}
\loigiai{Sắp xếp: $145$; $148$; $152$; $155$; $157$; $160$; $162$; $165$; $170$; $172$.

a) Sai: $\overline x=\dfrac{1586}{10}=158{,}6$.

b) Đúng: $R=172-145=27$.

c) Đúng: $Q_1=152$; $Q_3=165$; $\Delta_Q=13$.

d) Sai: $s^2=72{,}04>64$.}
\end{ex}
